% Part: first-order-logic
% Chapter: syntax-and-semantics
% Section: subformulas

\documentclass[../../../include/open-logic-section]{subfiles}

\begin{document}

\olfileid{fol}{syn}{sbf}

\olsection{\printtoken{P}{subformula}}

\begin{explain}
ఇచ్చిన \tetoken{సూత్రాలను}{!!{formula}s} ``ఏర్పరచే'' \tetoken{సూత్రం}{!!{formula}} గురించి మాట్లాడడం
తరచుగా ఉపయోగకరం. వాటిని దాని \emph{\tetoken{ఉపసూత్రాలు}{!!{subformula}s}} అంటాం. ప్రతి
\tetoken{సూత్రం}{!!{formula}} దానికే \tetoken{ఒక ఉపసూత్రంగా}{!!a{subformula}} లెక్కలోకి వస్తుంది; $!A$కు,
$!A$ స్వయంగా కాకుండా ఉన్న ఉపసూత్రాన్ని \emph{నిజ \tetoken{ఉపసూత్రం}{!!{subformula}}} అంటాం.
\end{explain}

\begin{defn}[తక్షణ \tetoken{ఉపసూత్రం}{!!^{subformula}}]
$!A$ \tetoken{ఒక సూత్రం}{!!a{formula}} అయితే, $!A$ యొక్క \emph{తక్షణ \tetoken{ఉపసూత్రాలు}{!!{subformula}s}}ను
కింది విధంగా ఆగమనాత్మకంగా నిర్వచిస్తాం:
\begin{enumerate}
\item పరమాణు \tetoken{సూత్రాలకు}{!!{formula}s} తక్షణ \tetoken{ఉపసూత్రాలు}{!!{subformula}s} ఉండవు.

\tagitem{prvNot}{\indcase{!A}{\lnot !B}{$\indfrm$కు ఉన్న ఏకైక తక్షణ
    \tetoken{ఉపసూత్రం}{!!{subformula}}~$!B$.}}{}

\item \indcase{!A}{(!B \ast !C)}{$\indfrm$ యొక్క తక్షణ
  \tetoken{ఉపసూత్రాలు}{!!{subformula}s} $!B$, $!C$ ($\ast$ ద్విస్థానిక సంయోజకాల్లో ఏదైనా).}

\tagitem{prvAll}{\indcase{!A}{\lforall[x][!B]}{$\indfrm$కు ఉన్న
    ఏకైక తక్షణ \tetoken{ఉపసూత్రం}{!!{subformula}}~$!B$.}}{}

\tagitem{prvEx}{\indcase{!A}{\lexists[x][!B]}{$\indfrm$కు ఉన్న
    ఏకైక తక్షణ \tetoken{ఉపసూత్రం}{!!{subformula}}~$!B$.}}{}
\end{enumerate}
\end{defn}

\begin{defn}[నిజ \tetoken{ఉపసూత్రం}{!!^{subformula}}]
$!A$ \tetoken{ఒక సూత్రం}{!!a{formula}} అయితే, $!A$ యొక్క \emph{నిజ \tetoken{ఉపసూత్రాలు}{!!{subformula}s}}ను
కింది విధంగా పునరావర్తకంగా నిర్వచిస్తాం:
\begin{enumerate}
\item పరమాణు \tetoken{సూత్రాలకు}{!!{formula}s} నిజ \tetoken{ఉపసూత్రాలు}{!!{subformula}s} ఉండవు.

\tagitem{prvNot}{\indcase{!A}{\lnot !B}{$\indfrm$ యొక్క నిజ
    \tetoken{ఉపసూత్రాలు}{!!{subformula}s}: $!B$, దానితో పాటు $!B$ యొక్క అన్ని నిజ
    \tetoken{ఉపసూత్రాలు}{!!{subformula}s}.}}{}

\item \indcase{!A}{(!B \ast !C)}{$\indfrm$ యొక్క నిజ
  \tetoken{ఉపసూత్రాలు}{!!{subformula}s}: $!B$, $!C$, వాటితో పాటు $!B$, $!C$ యొక్క అన్ని నిజ
  \tetoken{ఉపసూత్రాలు}{!!{subformula}s}.}

\tagitem{prvAll}{\indcase{!A}{\lforall[x][!B]}{$\indfrm$ యొక్క నిజ
    \tetoken{ఉపసూత్రాలు}{!!{subformula}s}: $!B$, దానితో పాటు $!B$ యొక్క అన్ని నిజ
    \tetoken{ఉపసూత్రాలు}{!!{subformula}s}.}}{}

\tagitem{prvEx}{\indcase{!A}{\lexists[x][!B]}{$\indfrm$ యొక్క నిజ
    \tetoken{ఉపసూత్రాలు}{!!{subformula}s}: $!B$, దానితో పాటు $!B$ యొక్క అన్ని నిజ
    \tetoken{ఉపసూత్రాలు}{!!{subformula}s}.}}{}
\end{enumerate}
\end{defn}

\begin{defn}[\tetoken{ఉపసూత్రం}{!!^{subformula}}]
$!A$ యొక్క \tetoken{ఉపసూత్రాలు}{!!{subformula}s}: $!A$ స్వయంగా, దానితో పాటు దాని అన్ని నిజ
\tetoken{ఉపసూత్రాలు}{!!{subformula}s}.
\end{defn}

\begin{explain}
తక్షణ \tetoken{ఉపసూత్రాలను}{!!{subformula}s}, నిజ \tetoken{ఉపసూత్రాలను}{!!{subformula}s} నిర్వచించిన తీరులోని
సూక్ష్మ భేదాన్ని గమనించండి. మొదటి సందర్భంలో $!A$కు సాధ్యమైన ప్రతి
ప్రతి $!A$ రూపం కోసం, దాని తక్షణ \tetoken{ఉపసూత్రాలను}{!!{subformula}s} నేరుగా నిర్వచించాం. ఇది
సందర్భాలవారీ స్పష్ట నిర్వచనం; ఈ సందర్భాలు \tetoken{సూత్రాలు}{!!{formula}s} సమితి యొక్క
ఆగమనాత్మక నిర్వచనాన్ని ప్రతిబింబిస్తాయి. రెండవ సందర్భంలో కూడా అన్ని
\tetoken{సూత్రాలు}{!!{formula}s} సమితిని నిర్వచించిన తీరునే ప్రతిబింబించాం; కానీ ప్రతి
సందర్భంలో చిన్న నిజ \tetoken{ఉపసూత్రాలు}{!!{subformula}s} అయిన \tetoken{సూత్రాలు}{!!{formula}s} $!B$, $!C$లను
మాత్రమే కాకుండా ఈ \tetoken{సూత్రాలను}{!!{formula}s} కూడా చేర్చాం. దీనివల్ల నిర్వచనం \emph{పునరావర్తకం}
అవుతుంది. సాధారణంగా ఆగమనాత్మకంగా నిర్వచించిన సమితిపై (మన సందర్భంలో
\tetoken{సూత్రాలపై}{!!{formula}s}) ఒక ప్రమేయపు నిర్వచనం, దాని సందర్భాల్లో అదే ప్రమేయాన్ని
తిరిగి ఉపయోగిస్తే పునరావర్తకమవుతుంది. అయితే అది సునిర్వచితం కావాలంటే,
మనం నిర్వచిస్తున్న ఆర్గ్యుమెంటు కంటే ``ముందు'' వచ్చే ఆర్గ్యుమెంట్లకు
ప్రమేయం విలువలను మాత్రమే వాడుతున్నామని నిర్ధారించాలి---మన సందర్భంలో
``నిజ \tetoken{ఉపసూత్రం}{!!{subformula}}''ను $(!B \ast !C)$కు నిర్వచించేటప్పుడు, ``ముందరి''
\tetoken{ఉపసూత్రాలు}{!!{subformula}s} అయిన \tetoken{సూత్రాలు}{!!{formula}s} $!B$, $!C$లనే వాడుతున్నాం.
\end{explain}

\begin{prop}
\ollabel{prop:subfrm-trans}
$!B$, $!A$కు ఉపసూత్రం, $!C$, $!B$కు ఉపసూత్రం అనుకుందాం. అప్పుడు
$!C$, $!A$కు ఉపసూత్రం. మరో మాటలో, ఉపసూత్ర సంబంధం సంక్రామకం.
\end{prop}

\begin{prob}
\olref[fol][syn][sbf]{prop:subfrm-trans}ను నిరూపించండి.
\end{prob}

\begin{prop}
\ollabel{prop:count-subfrms}
$!A$ అనేది $n$ సంయోజకాలు, పరిమాణీకరణికాలు గల సూత్రం అనుకుందాం.
అప్పుడు $!A$కు అత్యధికంగా $2n+1$ ఉపసూత్రాలు ఉంటాయి.
\end{prop}

\begin{prob}
\olref[fol][syn][sbf]{prop:count-subfrms}ను నిరూపించండి.
\end{prob}

\end{document}
